\chapter{Standards and Design Constraints}
This chapter documents the standards that the framework adheres to,
The constraints that have shaped its design and budgetary costs of the
Experiments; and concludes by mapping the project against the
Criteria for complex engineering problems and providing a brief summary.

\section{Compliance with the Standards}
Standards play out in three areas - the software stack, the hardware the
Framework executes upon, and the communication among components. There were
Alternative options available in each area and these alternatives are explained
Below.

\subsection{Software Standards}

Python 3.11 was chosen as the language to be used. This is because the LLM
Since Ecosystem is focussed on python, we also have hugging face libraries, vllm, and others.
Support for evaluation tools is in the form of python. While Java and c++ are faster
Both do not have the equivalent amount of pre-existing
Support to serve and load LLM models. Since the framework is a
Use a research prototype and not a production service, which must be real-time.
Ecosystem is better than raw speed. As such, surveys concerning LLM-based knowledge models are available.
Multi-agent systems discuss this design space in terms of workflow, infrastructure,
And open challenges~\cite{li2024survey} and therefore the decisions made
Below follow this structure.

vllm is built on top of the deep learning framework PyTorch 2.x. Pytorch is
The de-facto standard for LLM inference and directly supports qwen, Gemma and more.
Mistral families. While TensorFlow and JAX were evaluated alternatives;
However, TensorFlow's LLM serving story isn't as strong and JAX will require the whole pipeline.
On its own without any advantage.

Vllm serves the models via an OpenAI compatible api. This interface matters because
The framework will need to interact with closed APIs and serve locally stored models
With no modifications to the orchestration code. Given the requirements to handle
High throughput under the three agent debate load, PagedAttention provides vllm with
This capability. Additionally given FP8 support from PagedAttention matches that of
The Blackwell GPU. Ollama is simpler to setup and is slower than PagedAttention when handling concurrent requests. In contrast, llama.cpp runs well on CPU but does not match the GPU throughput needed for the debate. Tgi is a very similar alternative and vllm provided better ecosystem fit and an OpenAI compatible surface.

Cross-encoders can be used to re-rank passages returned by information retrieval algorithms. Sentence-transformers implements the cross-encoder that ranks the passages. In contrast, BM25 is an alternate ranking method however, like sentence-transformer, it only uses surface features to rank results. However, since both methods use the semantic features of a passage and a claim to score relevance between them, sentence transformers provide better overall relevance rankings. The web dashboard is served using FastAPI, and all result packages are created as JSON objects following rfc 8259, and are validated against a JSON schema. Codebase changes are tracked with git and latex (overleaf) generates this document.

\subsection{Hardware Standards}

Experiments are conducted on a single nvidia RTX PRO 6000 Blackwell GPU with 96 GB GDDR7 memory. Memory capacity and architecture were key considerations in making this selection. The development models (qwen3.5-9b, Gemma 4 12B, ministral-3-14b) fit into bf16. The final models (qwen3.6-27b, Gemma 4 26B, mistral small 3.2 24b) combine about 77 billion parameters; no single card holds this many parameters in fp16 so they run in FP8. Blackwell has native FP8 support; this is exactly why we chose it.

An initial alternative to consider was the a100 80gb. It has sufficient memory for development models and high-bandwidth HBM; however its FP8 path is less capable than what we have today, plus its architecture is two-generations old. Faster still is the h100 80gb; however it costs roughly double as much and offers none of what our framework requires. Finally, the L40S at 48 GB and the consumer RTX 5090 at 32 GB do not have enough memory to hold the full final model trio. Thus the RTX PRO 6000 sits in just the right place: there is enough memory to be useful; it has modern quantization; and it can be rented at a cost suitable for a student budget. To stage models prior to moving them to the GPU the host computer needs to have at least 64 GB of system RAM and nvme disk storage.

\subsection{Communication Standards}

All communication between the framework and its components occurs over http with tls. The model server presents an OpenAI-compatible chat completion api endpoint that is REST-based with JSON payloads. gRPC was another option for model serving; it has lower latencies in benchmark tests; but having an OpenAI-compatible surface allows us to make our framework portably usable across closed APIs and local vllm-served models without changing anything in our orchestration code.

Each source of evidence is accessed via its public REST interfaces: PubMed e-utilities, the arxiv api, the semantic scholar Graph API, and OpenAlex as a fallback. Each returns JSON or XML over HTTPS; our framework handles rate limits with retries capped according to those described in requirements. Web sockets are being saved for real-time updates in live trust charts in our dashboard.

\section{Design Constraints}

\subsection{Economic Constraint}


\subsection{Environmental Constraint}


\subsection{Ethical Constraint}


\subsection{Health and Safety Constraint}


\subsection{Social Constraint}


\subsection{Political Constraint}


\subsection{Sustainability}


\section{Cost Analysis}

The budget for this project consists of renting a GPU and for storing all data and for using the Research Services that do not require an additional charge to fine tune our results. Cost analysis is based almost entirely on the amount of time spent running our model on a rented GPU. The greatest expense (in terms of time) was inference time on the RTX PRO 6000 Blackwell which totaled approximately 300 hours during both development and testing. Hours were not distributed in proportion across each phase of the project. For example, Phase-0 and Phase-1 used about 30 total hours for setting up baseline models. Phase-2 also used about 60 total hours to develop the Claim Split and RAG models. Phase-3a and Phase-3b combined used about 140 hours for the Main Runs and Ablations. Phase-4 used about 40 total hours for conducting Human Study checks, while Phase-5 used about 30 hours for Final Runs and Packaging. These numbers show how Phase-3 had the most significant impact on costs. Table~\\ref{tab:budget} shows the major budget for the project and Table~\\ref{tab:budget-phase} details the number of GPU hours per phase.

\begin{center}
    \begin{table}[ht]
        \centering
        \caption{Primary budget for the project.}
        \renewcommand{\arraystretch}{1.6}
        \begin{tabular}{|p{0.32\textwidth}|p{0.18\textwidth}|p{0.18\textwidth}|p{0.18\textwidth}|}
            \hline
            Item                                         & Quantity  & Unit Cost (USD)   & Total (USD)       \\
            \hline
            RTX PRO 6000 Blackwell 96GB rental           & 300 hours & 0.70--1.90 / hour & 210--570          \\
            Retrieval APIs (PubMed, arXiv, S2, OpenAlex) & ---       & 0 (free tier)     & 0                 \\
            Datasets (GPQA, MMLU-Pro)                    & ---       & 0 (public)        & 0                 \\
            Storage (model weights, logs)                & 200 GB    & 0.05 / GB         & 10                \\
            \hline
            \textbf{Total}                               &           &                   & \textbf{220--580} \\
            \hline
        \end{tabular}
        \label{tab:budget}
    \end{table}
\end{center}

\begin{center}
    \begin{table}[ht]
        \centering
        \caption{GPU hours and cost split by phase (RTX PRO 6000, USD 0.70--1.90 per hour).}
        \renewcommand{\arraystretch}{1.6}
        \begin{tabular}{|p{0.32\textwidth}|p{0.18\textwidth}|p{0.18\textwidth}|}
            \hline
            Phase                                  & Hours        & Cost (USD)        \\
            \hline
            Phase-0 and Phase-1 (setup, baselines) & 30           & 21--57            \\
            Phase-2 (trust build)                  & 60           & 42--114           \\
            Phase-3a and Phase-3b (main runs)      & 140          & 98--266           \\
            Phase-4 (human eval)                   & 40           & 28--76            \\
            Phase-5 (writing, final checks)        & 30           & 21--57            \\
            \hline
            \textbf{Total}                         & \textbf{300} & \textbf{210--570} \\
            \hline
        \end{tabular}
        \label{tab:budget-phase}
    \end{table}
\end{center}

Costs for storage are estimated as \$10.00 based upon an initial cost of \$0.05/GB and a total of 200 GB. Logs and result package models will be approximately 180 GB in size while the remaining space will be used for the workspace. The direct costs of the retrieval API services from PubMed, arXiv, Semantic Scholar and OpenAlex are free due to their use of free tier service levels; however, the limits imposed may increase the time required to retrieve the data. As both GPQA and MMLU-Pro datasets are publically available, there were no additional costs associated with these datasets. Costs associated with printing (approximately \$15) and Internet connectivity were assumed to be fixed, therefore included in the overall estimate. Overleaf was provided to the team, thus eliminating any costs for this platform.

Additionally, an alternative budget was developed. Rented usage of an A100 80GB would be approximately \$0.68 - \$1.50/hour or 204 -- 450 dollars over the course of 300 hours. Although renting would produce significant savings over purchasing, the cost savings would be modest and come at a trade-off: reduced FP8 support from the Blackwell Path would necessitate tighter quantization or sequential loading, each of which would introduce engineering time and risk. Therefore, the RTX PRO 6000 was chosen because the Blackwell FP8 path would allow the pipeline to remain uncomplicated and avoid the engineering risks associated with those alternatives.

There is no revenue stream anticipated for this project. This project represents an academic deliverable produced by a university and a team with funding provided exclusively through grants awarded to the university and team. The only products expected from this effort include a written report and an open-source framework. An example of an alternate budget for comparison purposes is shown in Table~\\ref\{tab:budget-alt\}.


\begin{center}
    \begin{table}[ht]
        \centering
        \caption{Alternate budget using the A100 80GB.}
        \renewcommand{\arraystretch}{1.6}
        \begin{tabular}{|p{0.32\textwidth}|p{0.18\textwidth}|p{0.18\textwidth}|p{0.18\textwidth}|}
            \hline
            Item             & Quantity  & Unit Cost (USD)   & Total (USD)       \\
            \hline
            A100 80GB rental & 300 hours & 0.68--1.50 / hour & 204--450          \\
            Retrieval APIs   & ---       & 0 (free tier)     & 0                 \\
            Datasets         & ---       & 0 (public)        & 0                 \\
            Storage          & 200 GB    & 0.05 / GB         & 10                \\
            \hline
            \textbf{Total}   &           &                   & \textbf{214--460} \\
            \hline
        \end{tabular}
        \label{tab:budget-alt}
    \end{table}
\end{center}

A ten percent buffer (approximately \$25-\$55) is maintained to cover unexpected price increases and extended operating time in the event of a failed run. The buffer can be used to fund up to thirty additional hours of operation if either an ablation run fails or a maintenance call is required on a gate before the end of the plan period.


\section{Complex Engineering Problem}

\subsection{Complex Problem Solving}

The project meets the definition of a complex engineering problem by meeting the seven attributes defined in the Washington Accord criteria. Each attribute is mapped to subsections in Table~\ref{tab:p_solve}, which follows.

\begin{center}
    \begin{table}[ht]
        \centering
        \caption{Mapping with complex problem solving (Washington Accord P1--P7).}
        \renewcommand{\arraystretch}{1.6}
        \begin{tabular}{|p{0.11\textwidth}|p{0.11\textwidth}|p{0.11\textwidth}|p{0.11\textwidth}|p{0.11\textwidth}|p{0.11\textwidth}|p{0.11\textwidth}|}
            \hline
            P1                 & P2                                & P3                & P4                    & P5                         & P6                                & P7               \\
            Depth of Knowledge & Range of Conflicting Requirements & Depth of Analysis & Familiarity of Issues & Extent of Applicable Codes & Extent of Stakeholder Involvement & Inter-dependence \\
            \hline
            $\sqrt{}$          & $\sqrt{}$                         & $\sqrt{}$         & $\sqrt{}$             & $\sqrt{}$                  & $\times$                          & $\sqrt{}$        \\
            \hline
        \end{tabular}
        \label{tab:p_solve}
    \end{table}
\end{center}

\subsubsection{P1: Depth of Knowledge Required}

This project required knowledge of multiple areas including machine learning, natural language processing, retrieval-augmentation generation, and distributed model serving. Therefore, solving the project required specialized knowledge in these areas, such as the trust update formula through the vLLM serving layer. The coverage for P1 is based on the Washington Accord K1-K8, as detailed in Table~\ref{tab:knowledge} after P7.

\subsubsection{P2: Range of Conflicting Requirements}

P2 is covered. The project needed to meet multiple competing requirements. Accuracy competed with latency and cost. More rounds improved the signal but increased both the inference time and the GPU bill. Strong weighting would limit diversity due to trust calibration. The need to cover evidence competed with API rate limits. Therefore, the team needed to resolve these together. Round caps and clamping bounds were established, and verdicts are conservatively resolved. Thus, none of the competing needs dominate.

\subsubsection{P3: Depth of Analysis}

P3 is covered. There is no textbook solution to the problem. The team ran comparative experiments comparing several baseline systems: majority voting, MoA, iMAD, and DebUnc. Metrics were created to measure failure modes. Four types of metrics were developed: answer accuracy, consensus-collapse, minority-preservation, and evidence-calibration. An additional metric measuring the failure mode was the injection point for a fake wrong consensus; this allowed the team to measure collapse under controlled stress conditions. Therefore, the depth of test was sufficient to mark P3 as covered.

\subsubsection{P4: Familiarity of Issues}

P4 is covered. Multi-Agent LLM Debate is a novel research area. The most similar research that has been performed appears between 2024-2026; there is currently no standard method available for resolving issues related to sycophantic consensus in multi-agent LLM debates. At the beginning of the project, the issues were unfamiliar to the team. A complete review of current literature was conducted prior to designing any elements of the framework, as described in chapter 2. Therefore, this lack of familiarity is what makes P4 "covered".

\subsubsection{P5: Extent of Applicable Codes}

P5 is covered. General software standards are applicable. JSON (RFC 8259) and HTTP over TLS are applicable. Additionally, since the framework uses an open AI compatible API, the general API standard applies. However, as previously stated, there is currently no applicable code that addresses evidence-grounded trust in LLM debates. As such, this area is only partially standardized and the framework defines its own rules for the trust score and four possible verdicts. Therefore, this combination of existing codes and new conventions make P5 "covered."

\subsubsection{P6: Extent of Stakeholder Involvement}

P6 is not covered. The project does not have any external stakeholders. The only individuals involved are the supervisor and the examiners. There is no industry partner and no community group that identifies stakeholder needs or evaluates results as part of the development process. While users who will use future scientific QA tools may benefit from this project, they do not participate in either the design or testing phases of the project. The framework was designed and tested as an Academic FYDP (Final Year Design Project); therefore, the engagement is very narrow and only within the University context. For this reason P6 is considered "not covered", while the other P points are still considered covered.

\subsubsection{P7: Interdependence}

P7 is covered. The tight interdependencies between the Confidence Gate, Orchestrator, Claim Decomposer, Retrieval Subsystem, Trust Updater and Aggregator make it a candidate for P7. This means that if there is an error in the Retrieval Subsystem, this will be reflected in both Verdicts and ultimately the Final Answer through the Trust Score. Because of these interdependencies, P7 has been marked as "covered" with the addition of the Failure-Isolation Rule identified in Section 3.1.1.

Table~\ref{tab:knowledge} details the profile. The paragraphs below explain
each K for this project and why it is covered or not.

\begin{center}
    \begin{table}[ht]
        \centering
        \caption{Knowledge profile mapping for P1 (K1--K8).}
        \renewcommand{\arraystretch}{1.6}
        \begin{tabular}{|p{0.28\textwidth}|p{0.14\textwidth}|p{0.44\textwidth}|}
            \hline
            Knowledge                    & Applied   & Where used                                                        \\
            \hline
            K1: Natural sciences         & $\times$  & Not directly used; project is computing, not physics or chemistry \\
            K2: Mathematics              & $\sqrt{}$ & Softmax, clamping, renormalization, probability                   \\
            K3: Engineering fundamentals & $\sqrt{}$ & System design, state machines, failure handling                   \\
            K4: Specialist knowledge     & $\sqrt{}$ & LLMs, multi-agent debate, RAG, sycophancy                         \\
            K5: Engineering methods      & $\sqrt{}$ & Experimental methodology, baseline comparison                     \\
            K6: Computational methods    & $\sqrt{}$ & Python, PyTorch, vLLM, FastAPI                                    \\
            K7: Codes and practices      & $\sqrt{}$ & JSON, HTTP, TLS, version control                                  \\
            K8: Research and context     & $\sqrt{}$ & Ethical constraints and use in scientific QA                      \\
            \hline
        \end{tabular}
        \label{tab:knowledge}
    \end{table}
\end{center}

\paragraph{K1: Natural sciences} K1 is not covered. In K1, there isn't anything about natural science. There is no physics, no chemistry, and no biology — this is purely a computing project that has to do with how you use language models for retrieval. Because we're applying no part of natural science theory to the problem, we can mark K1 as "not addressed."

\paragraph{K2: Mathematics} K2 has been addressed. We apply softmax, clamping, and renormalization in our trust update; and we also have probability and statistical measures when using our metrics. Both were applied to our trust formula, and both were applied to our evaluation.

\paragraph{K3: Engineering fundamentals} K3 is covered. State machines (the system's design and its failure handling) are at the heart of what we've done here. Basic engineering design is used by requirements and the confidence gate along with the orchestration logic.

\paragraph{K4: Specialist knowledge} We've addressed K4. Our main focus of interest are LLMs (language models), multi-agent debate, RAG (Rebuttal Aggregation Graph), and sycophancy. This specialized way of thinking drives how we divide claims, how we retrieve information, and ultimately how we calculate trust.

\paragraph{K5: Engineering methods} K5 is covered. An experimental approach is being used with variable controls. Results are presented with confidence intervals and baselines are compared. All these studies were conducted in the same setup where injection was being tested.

\paragraph{K6: Computational methods} K6 is covered. Our main application tools are Python, PyTorch, vLLM (to serve the model), and FastAPI (to serve the dashboard). Sentence-transformers (to re-rank passages) and Git (for version control) complete the suite of applications.

\paragraph{K7: Codes and practices} K7 has been addressed. JSON under RFC 8259, HTTP with TLS, the OpenAI compatible API, and JSON Schema (as well as many others) were all applied as codes. Standard practices were followed regarding version control and coding style.

\paragraph{K8: Research and context} K8 has been addressed. Scientific QA (scientific questions answering) provides us with our ethical bounds on the output from AIs — literature reviews, gap analyses provide the context for our research direction.

The profile is broad rather than deep in one spot. The report applies K4 and K5
in every chapter, and K6 and K7 in the implementation, which shows the range the
problem demands.

\subsection{Engineering Activities}

Table~\ref{tab:e_act} maps the project to the five engineering activities
A1--A5, with rationales below.

\begin{center}
    \begin{table}[ht]
        \centering
        \caption{Mapping with complex engineering activities (A1--A5).}
        \renewcommand{\arraystretch}{1.6}

        \begin{tabular}{|p{0.16\textwidth}|p{0.16\textwidth}|p{0.16\textwidth}|p{0.16\textwidth}|p{0.16\textwidth}|}
            \hline
            A1                 & A2                   & A3         & A4                                       & A5          \\
            Range of resources & Level of Interaction & Innovation & Consequences for society and environment & Familiarity \\
            \hline
            $\sqrt{}$          & $\sqrt{}$            & $\times$   & $\times$                                 & $\sqrt{}$   \\
            \hline
        \end{tabular}
        \label{tab:e_act}
    \end{table}
\end{center}

\subsubsection{A1: Range of Resources}

A1 is covered. Three model family sources; Four Literature API's; Six team members; Rented Blackwell GPU. Coordinating these resources clearly has to be part of the work. The plan in Chapter 3 assigns each task to a specific team member and tracks weeks/deliverables. That is why A1 is labeled as covered.

\subsubsection{A2: Level of Interaction}

A2 is covered. The system makes heavy use of the outside world. Every round of debating calls three external model API endpoints and three Literature API endpoints, and there is communication to/from the rented cloud GPU. Team communication occurs via the shared repository and weekly meeting. High degree of interactivity = why A2 is labeled as covered.

\subsubsection{A3: Innovation}

A3 is not covered. The framework combines pre-existing concepts; multi-agent debate/RAG with trust-weighting/reranking. It creates no new device/new theory. The innovation is a bound trust-score that updates based upon evidence; useful, but still within an incremental level of innovation. No patent-able/break-through innovation exists. That is why A3 is labeled as "not covered", although it would represent a first time combination of concepts for this project.

\subsubsection{A4: Consequences for Society and Environment}

A4 is not covered. The work represents a lab prototype, NOT a production-deployed product. Any social-impact will be limited to providing more-reliable research results from a test-environment, rather than changing daily life/social habits, or supporting a large user-base. Similarly, environmental-impact will be minimal. Test uses approximately 300 GPU-hours on public-cloud using standard configurations/no-fine tuning - therefore energy usage will remain very-low. No large-scale deployment exists; therefore, broad societal-consequences do not exist. That is why A4 is labeled as "not covered".

\subsubsection{A5: Familiarity}

A5 is covered. Both multi-agent debate/evidence verification were uncharted territory for our team. Our literature-review phase and our baseline study that used a step-by-step implementation plan were both structured around the lack of prior experience with those areas. Only after we had studied each component sufficiently could we begin implementing each one. Lack of prior-experience with many components = why A5 is labeled as "covered".

\section{Summary}
The chapter defined the requirements for the software-layer, hardware-layer, and communication-layers and documented the economic, environmental, ethics-social constraints that have impacted design decisions. Budget-analysis demonstrates that the project can fit within a students' budget utilizing a rented RTX PRO 6000 Blackwell GPU with a phased split of costs over the 300 hours plus a 10\% buffer. The complex engineering mapping demonstrates that six out of seven P attributes are covered with only P6 being non-covered due to little to no input from external stakeholders. For knowledge, K2-K8 are covered while K1 is not. For activities, A1-A2 and A5 are covered while A3 and A4 are not.
